%%%%%%%%%%  scan idx 674  |  folio 661  %%%%%%%%%%
\SGAPrintPage{661}%
\sgaStableTarget{sga6.e13.cor.6-14}\textbf{Corollary 6.14. (G. Hermann).} \emph{For every $N\ge0$, there exists a polynomial $R_N(m)$ such that, for every field $k_0$ and every ideal $I$ of the polynomial ring $k_0[T_1,\ldots,T_N]$, generated by elements of degree $\le m$, the radical $\sqrt I$ of $I$ is generated by elements of degree $\le R_N(m)$.}

Indeed, let $k$ be the algebraic closure of $k_0$. Introducing an auxiliary indeterminate, consider $Y$ (resp. $X$), the subscheme of $P=\mathbb P^N_k$ defined by the homogenized ideal of $I$ (resp. of $\sqrt I$). In the Hilbert polynomial $\chi({O}_Y(n))=\sum_{i=0}^{N}b_i\binom{n+i}{i}$, we have $|b_i|\le c_i$ with $c_i=A_{N-i}(0,\ldots,0,1;m)$ by 6.13.

Let $X=\bigcup X^q$ be the decomposition according to dimension. Then $\text{degree}(X^q)\le e_q$ with $e_q=P_{N-q}(c_q,\ldots,c_N)$. Indeed, after cutting $X$ (resp. $Y$) by a ``general'' linear space of codimension $q$, we may suppose that $q=0$. Then clearly, $\text{degree}(X^q)=h^0({O}_{X^q})$ and $h^0({O}_{X^q})\le h^0({O}_Y)$, and by 6.4 (iii), $h^0({O}_Y)\le e_q$.

Consequently, ${O}_{X^q}$ is a $(0,\ldots,0,e_q)$-sheaf by 6.10. It therefore follows from the injection ${O}_X\to \prod{O}_{X^q}$ that ${O}_X$ is an $(e_0,\ldots,e_N)$-sheaf. Thus, in the Hilbert polynomial $\chi({O}_X(n))=\sum_{i=0}^{N}a_i\binom{n+i}{i}$, we have $|a_i|\le f_i$ with $f_i=A_{N-i}(e_i,\ldots,e_N;0)$ by 6.7. Finally, the ideal $\sqrt I$ of $X$ is $R_N(m)$-regular with $R_N(m)=P_N(f_0,\ldots,f_N)$; whence the assertion by 1.3.

%%%%%%%%%%  scan idx 675  |  folio 662  %%%%%%%%%%
\SGAPrintPage{662}%
\sgaStableTarget{sga6.e13.sec.7}\subsection*{7. Appendix: The Hodge index theorem}

Here we recover and complete some earlier results by a direct argument.

\sgaStableTarget{sga6.e13.thm.7-1}\textbf{Theorem 7.1.} \emph{Let $X$ be an irreducible surface, proper over an algebraically closed field. Suppose there exists an invertible sheaf $H$ on $X$ such that $\langle c_1(H)^2\rangle>0$ (for example, $H$ ample). Then an invertible sheaf $L$ on $X$ such that}
\[
   \langle c_1(L)\cdot c_1(H)\rangle=0
   \quad\hbox{\uline{and}}\quad
   \langle c_1(L)^2\rangle\geq 0,
\]
\emph{is numerically equivalent to $0$. (Consequently, the intersection form on $A^1_{\mathrm{num}}(X)_{\Q}$ ($=\bigl[\Pic(X)/\Pic^\tau(X)\bigr]\otimes\Q$, by 4.6) is nondegenerate of type $(1,\rho-1)$).}

\sgaStableTarget{sga6.e13.para.7-1-1}\textbf{7.1.1.} First suppose: $X$ reduced (hence integral); $H$ ample; $\langle c_1(L)\cdot c_1(H)\rangle=0$; and $\langle c_1(L)^2\rangle>0$. Arguing by contradiction, take $m>0$ such that $H_1=L\otimes H^{\otimes m}$ is ample, and take $n$ such that, putting $N=L^{\otimes n}\otimes H^{-1}$, we have that $\langle c_1(N)\cdot c_1(H_1)\rangle=n\langle c_1(L)^2\rangle-m\langle c_1(H)^2\rangle$ is $>0$. Then we also have that $\langle c_1(N)\cdot c_1(H)\rangle=-\langle c_1(H)^2\rangle$ is $<0$, and that $\langle c_1(N)^2\rangle=n^2\langle c_1(L)^2\rangle+\langle c_1(H)^2\rangle$ is $>0$, which contradicts the following lemma.

\sgaStableTarget{sga6.e13.lem.7-1-2}\textbf{Lemma 7.1.2.} \emph{Let $X$ be an integral surface, projective over an algebraically closed field. For an invertible sheaf $N$ on $X$ such that $\langle c_1(N)^2\rangle>0$, the following conditions are equivalent}:
\begin{enumerate}[label=(\roman*)]
\item \emph{For every ample sheaf $H$, we have $\langle c_1(N)\cdot c_1(H)\rangle>0$.}
\item[(i')] \emph{There exists an ample sheaf $H$ such that we have $\langle c_1(N)\cdot c_1(H)\rangle>0$.}
\item \emph{There exists an integer $n>0$ such that we have $H^0(X,N^{\otimes n})\ne0$.}
\end{enumerate}

%%%%%%%%%%  scan idx 676  |  folio 663  %%%%%%%%%%
\SGAPrintPage{663}%
Indeed, (ii) implies that there exists an effective divisor $D$ such that $N^{\otimes n}\simeq{O}_X(D)$, and $\langle c_1(N)^2\rangle>0$ implies that $D>0$; whence (i).

To see that (i') implies (ii), we may clearly suppose that $H={O}_X(1)$ is very ample; then there exists a closed integral, locally principal curve $Y$ such that $H={O}_X(Y)$. By Riemann's theorem, (i') implies that $H^1\bigl(N_Y^{\otimes n}(p)\bigr)=0$ for $p\geq0$ and $n\geq n_0$ with $n_0=2h^1({O}_Y)-1$. Since $H^2\bigl(N^{\otimes n}(p)\bigr)=0$ for $p\gg0$, the exact sequence
\[
   H^1\bigl(N_Y^{\otimes n}(p)\bigr)
   \longrightarrow H^2\bigl(N^{\otimes n}(p-1)\bigr)
   \longrightarrow H^2\bigl(N^{\otimes n}(p)\bigr)
   \longrightarrow0
\]
gives $H^2(N^{\otimes n})=0$ for $n\geq n_0$. Thus, for $n\gg0$, we have
\[
   h^0(N^{\otimes n})
   =\chi(N^{\otimes n})+h^1(N^{\otimes n})
   \geq\chi(N^{\otimes n})
   =1/2\langle c_1(N)^2\rangle n^2+\ldots>0,
\]
which completes the proof of the lemma.

\sgaStableTarget{sga6.e13.para.7-1-3}\textbf{7.1.3.} Now suppose in 7.1 that there exists an invertible sheaf $M$ on $X$ such that $\langle c_1(L)\cdot c_1(M)\rangle\ne0$, and argue by contradiction. Replace $X$ by the underlying reduced scheme of a Chow presentation of $X$, which is allowed by the projection formula X 4.4.4, and let $G$ be ample on $X$. Take $m>0$ such that $H_1=G^{\otimes m}\otimes M$ is ample, and such that $\langle c_1(L)\cdot c_1(H_1)\rangle\ne0$; take $p,q$ ($q\ne0$) such that, putting $L_1=L^{\otimes p}\otimes H^{\otimes q}$, we have that $\langle c_1(L_1)\cdot c_1(H_1)\rangle=p\langle c_1(L)\cdot c_1(H_1)\rangle+q\langle c_1(H)\cdot c_1(H_1)\rangle$ is $0$; then we also have that $\langle c_1(L_1)^2\rangle=p^2\langle c_1(L)^2\rangle+q^2\langle c_1(H)^2\rangle$ is $>0$; this returns us to case 7.1.1, with $L=L_1$ and $H=H_1$.

\sgaStableTarget{sga6.e13.para.7-1-4}\textbf{7.1.4.} Finally, consider the general case of 7.1. Still arguing by contradiction, suppose that $L$ is not numerically equivalent to $0$; in addition, replace $X$ by the normalization of $X_{\red}$, so that $X$ has
%%%%%%%%%%  scan idx 677  |  folio 664  %%%%%%%%%%
\SGAPrintPage{664}%
only finitely many singular points, a reduction justified by the projection formula X 4.4.4 and by 4.5 (ii). There therefore exists a closed integral curve $Y$ on $X$ such that $\langle c_1(L)\cdot Y\rangle\ne0$, and $Y$ is locally principal outside a finite set $Z$. Let $f:X'\to X$ be the blowup of the ideal $I$ of ${O}_X$. The ideal $M=I\cdot{O}_{X'}$ is invertible, and defines a subscheme of $X'$ which is the union of two subschemes $Y'$, $Y''$ such that $f(Y')=Y$ and $f(Y'')\subset Z$. By the projection formula X 4.4.4, we have that $\langle c_1(f^*L)\cdot c_1(M)\rangle=-\langle c_1(f^*L)\cdot Y'\rangle-\langle c_1(f^*L)\cdot Y''\rangle$ equals $-\langle c_1(L)\cdot Y\rangle-0\ne0$, and we may replace $X,L,H$ by $X',f^*L,f^*H$; we are therefore reduced to the preceding case, and the proof of 7.1 is complete.

\sgaStableTarget{sga6.e13.rem.7-2}\textbf{Remark 7.2.} It is unknown whether an irreducible proper surface (over an algebraically closed field), which possesses an invertible sheaf $H$ such that $\langle c_1(H)^2\rangle>0$, is necessarily projective. It is also unknown whether, for every irreducible proper surface, the intersection form on $A^1_{\mathrm{num}}(X)_{\Q}$ is nondegenerate; it follows from 7.1 that, when it is degenerate, it is necessarily negative (nondefinite).

\sgaStableTarget{sga6.e13.cor.7-3}\textbf{Corollary 7.3.} \emph{Under the conditions of 7.1, let $M$ be an invertible sheaf on $X$ which is not numerically proportional to $H$, and let $d=\langle c_1(M),c_1(H)\rangle$ and $h=\langle c_1(H)^2\rangle$. Then, $\langle c_1(M)^2\rangle<d^2/h$.}

Indeed, put $L=M^{\otimes h}\otimes H^{\otimes -d}$. Since $\langle c_1(L)\cdot c_1(H)\rangle=0$, by 7.1, $\langle c_1(L)^2\rangle<0$; whence the assertion.

%%%%%%%%%%  scan idx 678  |  folio 665  %%%%%%%%%%
\SGAPrintPage{665}%
\sgaStableTarget{sga6.e13.cor.7-4}\textbf{Corollary 7.4.} \emph{Let $X$ be an irreducible scheme of dimension $r\geq2$, projective over an algebraically closed field, and let $H$ be an ample sheaf on $X$. Then}:
\begin{enumerate}[label=(\roman*)]
\item \emph{For an invertible sheaf $L$ on $X$ to be numerically equivalent to $0$, it is (necessary and) sufficient that}
\[
   \langle c_1(L)\cdot c_1(H)^{r-1}\rangle=0
   \quad\hbox{\uline{and}}\quad
   \langle c_1(L)^2\cdot c_1(H)^{r-2}\rangle\geq0.
\]
\item \emph{For every invertible sheaf $M$ on $X$ which is not numerically proportional to $H$, we have}
\[
   \langle c_1(M)^2\cdot c_1(H)^{r-2}\rangle
   <
   \langle c_1(M)\cdot c_1(H)^{r-1}\rangle^2/\langle c_1(H)^r\rangle.
\]
\end{enumerate}

Indeed, (ii) follows immediately from (i), by the argument of 7.3. To establish (i), let $Y$ be a closed integral curve in $X$; we wish to prove that $\langle c_1(L)\cdot Y\rangle=0$. Embed $X$ in $\mathbb P^N_k$ by means of $H^{\otimes n}$ ($n>0$). Let $I=\bigoplus I_\nu$ be the homogeneous ideal of $Y$, and let $\nu_0$ be an integer such that $I_{\nu_0}$ generates $I_\nu$ for $\nu\geq\nu_0$. Then the hypersurfaces $Z$ of degree $\nu_0$ that contain $Y$ have no common points outside $Y$. Arguing by induction on $r=\dim(X)$, we find, using Bertini's theorem, $r-2$ hypersurfaces $Z_1,\ldots,Z_{r-2}$ of degree $\nu_0$ that contain $Y$, such that $X'=Z_1\cap\ldots\cap Z_{r-2}\cap X$ is irreducible. Since $\langle c_1(L)\cdot Y\rangle=\langle c_1(L_{X'})\cdot Y\rangle$, it suffices to prove that $L_{X'}$ is numerically equivalent to $0$ on $X'$. Now, by the projection formula, we have $\langle c_1(L_{X'})\cdot c_1(H_{X'})\rangle=0$ and $\langle c_1(L_{X'})^2\rangle\geq0$; thus the assertion follows from 7.1.

\sgaStableTarget{sga6.e13.cor.7-5}\textbf{Corollary 7.5.}
\emph{Let $k$ be an algebraically closed field, $Y$ a projective $k$-scheme equipped with an ample sheaf $H$, and $X$ the zero scheme of a section $\varphi$ of $H$. Suppose that, for every irreducible component $Y^i$ of $Y$,
%%%%%%%%%%  scan idx 679  |  folio 666  %%%%%%%%%%
\SGAPrintPage{666}%
the irreducible components of $X\cap Y^i=V(\varphi|Y^i)$ have dimension $\geq2$. Then the morphism $A^1_{\mathrm{num}}(Y)\longrightarrow A^1_{\mathrm{num}}(X)$ is injective.}

Indeed, since we may clearly suppose that $Y$ is integral, the assertion follows from 7.4 (i), by the projection formula.

\sgaStableTarget{sga6.e13.bibliography}\subsection*{BIBLIOGRAPHY}
\begin{enumerate}[label={[\arabic*]}]
\item Abhyankar S. : Resolution of singularities of arithmetical surfaces, in : Arithmetical Algebraic Geometry, Purdue University, Harper-Row, New York 1965.
\item Mumford D : Pathologies III, Amer. J. Math 1967 p. 94-104.
\item Grothendieck A. and Dieudonn\'e J. EGA, El\'ements de G\'eom\'etrie Alg\'ebrique, Pub. Math. Inst des Hautes Etudes Sci., 1960.
\item Grothendieck A. : FGA, Fondements de la g\'eom\'etrie alg\'ebrique, Extracts from the S\'eminaire Bourbaki 1957-1962.
\item Grothendieck A. et alii, SGA, S\'eminaire de g\'eom\'etrie alg\'ebrique, 1960-66, I.H.E.S.
\item Jussila O. : X, Formalisme des intersections sur les sch\'emas alg\'ebriques propres, Expos\'e X in this seminar.
\item Raynaud M. : XII, Un th\'eor\`eme de repr\'esentabilit\'e relative sur le foncteur de Picard, Expos\'e XII in this seminar.
\end{enumerate}



% Exposé XIV local notation
\providecommand{\Kb}{K^\bullet}
\providecommand{\Db}{D}
\providecommand{\Gtop}{\operatorname{Gr}_{\mathrm{top}}}
\providecommand{\Ftop}{\operatorname{Fil}_{\mathrm{top}}}


